\section{Related work}\label{sec:related}

\paragraph{Heuristic extremal maps and ensembles.}
Most computational work on extreme partisan maps is heuristic. Recombination and sequential Monte Carlo samplers \citep{deford2021,mccartan2023} generate ensembles; short-burst hill climbing \citep{cannon2023} and column generation \citep{gurnee2021} optimise chosen objectives; multilevel schemes \citep{swamy2023} approximate Pareto fronts; Ising-model and integer-programming-plus-local-search heuristics target maximally tilted assignments \citep{okamoto2021} and majority-minority districts \citep{brous2025}, and a mixed-integer model combines district drawing with targeted campaigning \citep{deshpande2023}; and recent studies map how large the achievable partisan swings are \citep{goedert2024} and how durable optimised majorities are under electoral change \citep{palmer2026}. The public Algorithmic Redistricting Atlas publishes short-burst ``Max'' (51\%) and ``Safe Max'' (55\%) maps and states explicitly that they carry no optimality guarantee \citep{atlas2026}. Ensemble analyses place enacted plans within a neutral baseline (for New Hampshire, \citealp{mcwhorter2026}); Section~\ref{sec:enacted} places them instead relative to certified extremes. These methods produce feasible plans---lower bounds on our frontier---but not proofs that nothing better exists. We use a county-aware recombination hill climber only as a source of verified incumbents.

\paragraph{Exact integer programming for districting.}
Integer programming for districting goes back to \citet{hess1965}. The key modelling advance is exact contiguity: \citet{validi2022} give branch-and-cut formulations with lazily separated connectivity cuts, \citet{miyazawa2021} study formulations for balanced connected partitions, and \citet{jolly2026} catalogue \emph{inexact} contiguity constraints that can exclude valid maps---which would invalidate any claimed upper bound. \citet{shahmizad2025} compute the exact minimum number of county splits for every state. Those works optimise compactness or splits; none optimises a partisan seat-count objective with a proof.

\paragraph{Bounds for representation objectives.}
Two lines are close to ours. \citet{duchin2019} derive rigorous \emph{geography-free} infeasibility bounds for the number of districts a party can win in Massachusetts by sorting towns or precincts by partisan margin per capita (with a pseudo-polynomial knapsack dynamic programme where sorting is inconclusive), for districts built from arbitrary collections of units without any spatial requirement; our geography-free baseline (Section~\ref{sec:negative}) is the same idea extended to the whole margin spectrum and to $k$-way allocations, in the spirit of the theory of optimal gerrymandering with unconstrained geography \citep{owen1988,friedman2008}, which is also studied as an optimisation problem in its own right \citep{benade2021} and, with uncertainty about district-level outcomes, in a stylised continuous model \citep{lagarde2024}. On the theory side, maximising the seats of a party is NP-complete on planar graphs and admits constant-factor approximations there \citep{dippel2024}; our exact and certified results are computational, for the graphs that arise from real precincts. \citet{fravel2026} prove dual bounds for non-convex expected-representation objectives with exact contiguity, but for whole-county instances only, with a 10\% population tolerance and at most seven districts. Because forcing counties whole \emph{restricts} the plan class, a county-level optimum is not an upper bound for maps with finer resolution; the relaxations of Section~\ref{sec:relax} are valid at precinct resolution by construction.

\paragraph{Adaptive refinement.}
Our algorithm is an instance of counterexample-guided abstraction refinement \citep{clarke2003} and is analogous to dynamic discretization discovery in transportation \citep{boland2017,marshall2021}: solve a small relaxation whose optimum is a valid bound, refine only where the relaxed solution is spurious, and stop when the answer is proved or an exact solution appears. What is specific to districting is the abstraction: cells with exact-integer aggregate allocations, fractional-knapsack hull cuts and quotient-graph connectivity, together with a subdivision budget that makes coarse abstractions informative.
